\documentclass{llncs}

\usepackage[T1]{fontenc}
\usepackage[utf8]{inputenc}
\usepackage[spanish,es-nodecimaldot,es-tabla]{babel}
\usepackage{graphicx}
\usepackage{float}
\usepackage{hyperref}
\usepackage{url}

\renewcommand\UrlFont{\color{blue}\rmfamily}
\urlstyle{rm}

\raggedbottom

\begin{document}

\title{NeuroRisk: Predicción Temprana de Alteraciones en el Neurodesarrollo de Neonatos Prematuros}

\author{Juan Pablo Buriticá Ruiz \and
Mateo Hernández Gualdrón \and
Martín Alejandro Quintero Rincón \and
Andrés Felipe Romero Grajales \and
María Paula Saa Beltrán}

\institute{Universidad de los Andes, Maestría en Inteligencia Artificial\\
Producto de Datos, Micro-proyecto, Entrega 3\\
Bogotá, Colombia, septiembre de 2026}

\maketitle

% =====================================================================
\section{Resumen del problema}

\textbf{Contexto.} La prematurez y el bajo peso al nacer son factores de riesgo
establecidos para alteraciones del neurodesarrollo en los dominios cognitivo, de
lenguaje, perceptual y motor. Estas alteraciones se consolidan en los primeros
años de vida, lo que abre una ventana de intervención temprana si se identifica
a tiempo qué neonatos tienen mayor riesgo.

\textbf{Pregunta de negocio y alcance.} ¿Cuáles neonatos prematuros tienen mayor
riesgo de presentar una alteración del neurodesarrollo, de modo que el equipo
clínico pueda priorizar, desde el momento del alta, el seguimiento y la
intervención temprana en los pacientes que más lo necesitan? La predicción se hace
\emph{al alta}, cuando ya se conocen las duraciones de hospitalización y de
oxígeno. El prototipo incluye DVC, MLflow, una API en FastAPI, un tablero en
Streamlit y despliegue con Docker Compose.

\textbf{Datos.} Cohorte retrospectiva de 89 neonatos prematuros del Hospital
Ghaem (Mashhad, Irán), 2016 a 2020, con 53 variables maternas, neonatales, de
intervenciones y de evaluación Bayley \cite{darabi2024,bayley}. El target
\texttt{neurodev\_alteration} vale 1 si al menos uno de los cinco dominios
Bayley es anormal (52.8\% de casos positivos). Los desenlaces se midieron en una
única ronda de evaluación en 2021 \cite{faramarzi2023}, con edades entre 4 meses
y 3.3 años; por eso \texttt{Age} y \texttt{correctedage} se excluyen como
predictores (son la edad al desenlace, no información disponible al alta). Al
provenir de un solo hospital, los resultados no deben asumirse transferibles a
otros contextos sin validación externa.

\textbf{Cambios respecto a la Entrega 2.} (i) Validación cruzada repetida en
lugar del holdout de 18 filas, con tuning y un método explícito de selección de
variables; (ii) modelo final de regresión logística elástica con odds ratios
reportados; (iii) formulario con \texttt{duration.hopitalization} y
\texttt{duration.O2}, explicación local del riesgo y listado de priorización;
(iv) despliegue con Docker Compose y MLflow compartido en EC2. Se corrige además
el argumento del cambio de target: la prevalencia de sepsis (5/89) es la real de
la condición, no un defecto del dataset; el problema es que con un split 80/20
estratificado el conjunto de prueba tendría un solo caso positivo, lo que
impide estimar recall o ROC-AUC. Por eso se usa un target compuesto y la sepsis
queda como predictor.

% =====================================================================
\section{Modelos desarrollados y su evaluación}

\subsection{Preprocesamiento}

El pipeline de features (\texttt{build\_features.py}, registrado como stage de
DVC) codifica variables binarias y categóricas, corrige un error
de captura conocido (\texttt{congenital.anomaly = 34.0}, que duplica la edad
gestacional de la misma fila, se trata como nulo) y combina
\texttt{surfactant} y \texttt{aggressive.ventilation} (coinciden en 95.5\% de
los casos) en un solo indicador. La imputación y el escalado
(\texttt{RobustScaler}) se ajustan dentro de un \texttt{sklearn.Pipeline} solo
sobre los datos de entrenamiento de cada fold, para evitar fuga de información.

\subsection{Selección de algoritmos}

Se compararon tres familias que cubren supuestos distintos sobre la relación
entre variables y riesgo:
\emph{regresión logística}, lineal e interpretable, adecuada cuando hay pocas
observaciones y relaciones aproximadamente monótonas como las observadas en el
EDA;
\emph{Random Forest}, un ensamble por bagging que captura no linealidades e
interacciones y es robusto con poco ajuste;
y \emph{XGBoost}, un ensamble por boosting que suele ser el referente en datos
tabulares. Como variaciones adicionales se probaron pesos de clase, SVM, árboles
extremadamente aleatorizados, selección de variables con L1 y sobremuestreo con
SMOTENC. Todo quedó registrado en MLflow (21 runs en el experimento
\texttt{optimizacion\_modelo\_final}, Figura~\ref{fig:mlflow}).

\subsection{Selección de variables}

Se construyó un ranking combinando cuatro criterios: poder univariado,
estabilidad de la selección L1 en 200 submuestras, e importancia por permutación
con dos modelos distintos (\texttt{docs/ranking\_variables.csv}). Tres variables
dominan y son las únicas con estabilidad L1 alta: \texttt{BirthWeight} (a menor
peso, mayor riesgo), \texttt{duration.hopitalization} y \texttt{duration.O2} (a
más días, mayor riesgo). Después el ranking se aplana: siguen
\texttt{SepsisnegativeCulture}, \texttt{PregnancyAge}, \texttt{IUGR},
\texttt{icter} y \texttt{BPD}. Varias complicaciones son poco frecuentes
(\texttt{meningitis} no tiene casos positivos y otras seis variables binarias
tienen cinco o menos). Esto refleja cómo se comporta el fenómeno en una cohorte
de 89 neonatos, no un defecto del dato; su consecuencia es que el efecto de
esas variables se estima con mucha incertidumbre.

\subsection{Validación, tuning y criterio de selección}

Se usó validación cruzada estratificada repetida (5 folds por 3 repeticiones)
sobre las 89 filas. Maximizar recall con umbral 0.5 resultó engañoso: los
modelos muy regularizados devolvían probabilidades cercanas a la prevalencia y
marcaban casi todo como anormal (especificidad de 0.17). Por eso el criterio
principal es el mayor recall alcanzable exigiendo una precisión mínima de 0.65,
junto con un control de la dispersión de las probabilidades predichas.

\begin{table}[H]
\centering
\caption{Comparación de modelos con validación cruzada estratificada repetida (5$\times$3).}
\label{tab:modelos}
\begin{tabular}{|l|c|c|c|}
\hline
Modelo & ROC-AUC & Recall (0.5) & Recall (prec. $\geq$ 0.65) \\
\hline
Clasificador trivial (siempre anormal) & 0.500 & 1.000 & 0.000 \\
Baseline regresión logística & 0.602 & 0.518 & 0.383 \\
Baseline Random Forest & 0.667 & 0.646 & 0.582 \\
Baseline XGBoost & 0.590 & 0.610 & 0.418 \\
Random Forest ajustado & 0.723 & 0.637 & 0.766 \\
\textbf{Regresión logística elástica (final)} & \textbf{0.691} & \textbf{0.623} & \textbf{0.638} \\
\hline
\end{tabular}
\end{table}

\begin{figure}[H]
\centering
\includegraphics[width=\textwidth]{mlflow_runs.png}
\caption{Vista de runs del experimento \texttt{optimizacion\_modelo\_final} en
el servidor compartido de MLflow en EC2: métricas de validación cruzada (media y
desviación estándar) por configuración evaluada.}
\label{fig:mlflow}
\end{figure}

\subsection{Modelo final}

Se eligió una regresión logística con penalización elástica ($C=0.01$,
$l_1$\_ratio $=0.05$), reentrenada con las 89 filas y exportada como un único
\texttt{.joblib} en MLflow. Siete modelos quedaron estadísticamente empatados
(el error estándar del recall ronda 0.18), y entre ellos se prefirió el más
parsimonioso e interpretable. El Random Forest ajustado obtiene un recall
mayor con la restricción de precisión (0.766 frente a 0.638), pero la
diferencia está dentro del ruido de la validación y no ofrece coeficientes
interpretables para el equipo clínico; queda documentado como alternativa.

\textbf{Odds ratios.} Ninguna asociación es firme con esta muestra. Solo
\texttt{sepsis} (OR 2.10, intervalo 1.01 a 3.28, con cinco casos),
\texttt{PROM} y \texttt{pneumothorax} tienen intervalos que no incluyen 1. Los
efectos de las variables continuas van en la dirección clínicamente esperada
(menor peso al nacer y estancias más largas aumentan el riesgo), pero con
intervalos anchos. La tabla completa está en
\texttt{notebooks/model\_optimization.ipynb}.

% =====================================================================
\section{Tablero desarrollado y funcionalidad}

El tablero (\texttt{dashboard/app.py}) consume la API
(\texttt{api/main.py}), que carga el modelo desde MLflow y expone
\texttt{GET /health} y \texttt{POST /predict}. Responde a la pregunta de
negocio con tres funcionalidades:

\textbf{Calculadora de riesgo al alta.} El formulario pide las variables que más
pesan en el modelo, incluidas \texttt{duration.hopitalization} y
\texttt{duration.O2}, que en la Entrega 2 se completaban con valores por
defecto de la cohorte. Esa era una limitación importante: el puntaje quedaba
dominado por valores que no eran del paciente. Pedir solo \texttt{BirthWeight}
y las dos duraciones ya sube el ROC-AUC en validación cruzada de 0.65 a 0.72.
El resultado muestra probabilidad, puntaje de 0 a 100, categoría de riesgo y
recomendación clínica.

\textbf{Explicación local.} Aprovechando la interpretabilidad del modelo lineal,
para cada paciente se multiplica el valor escalado de cada variable por su
coeficiente y se grafican las contribuciones ordenadas por magnitud: en rojo lo
que sube el riesgo y en verde lo que lo baja (Figura~\ref{fig:tablero}). Como el
escalado está centrado en la mediana, cada contribución se lee respecto a un
paciente típico de la cohorte.

\textbf{Priorización.} Una pestaña dedicada, el panel de triaje, consolida los
neonatos evaluados durante la sesión clínica y los ordena de mayor a menor
riesgo estimado, con filtros por categoría y exportación a CSV, para decidir a
quién priorizar en el seguimiento (Figura~\ref{fig:triage}). La pestaña de
exploración conserva los KPIs descriptivos de la cohorte.

\begin{figure}[H]
\centering
\includegraphics[width=\textwidth]{dashboard_contribuciones.png}
\caption{Calculadora con la explicación local: variables que empujan el riesgo
del paciente hacia arriba (rojo) y hacia abajo (verde).}
\label{fig:tablero}
\end{figure}

\begin{figure}[H]
\centering
\includegraphics[width=\textwidth]{dashboard_triage.png}
\caption{Panel de triaje: neonatos evaluados en la sesión, ordenados por
riesgo estimado.}
\label{fig:triage}
\end{figure}

\textbf{Despliegue.} Docker Compose (\texttt{deploy/docker-compose.yml}) levanta
dos servicios en una red interna: la API en el puerto 8000 y el tablero en el
8501, que llama a la API por nombre de servicio (\texttt{http://api:8000}). Se
montó además un volumen con el \texttt{.joblib} como respaldo, para que el
despliegue funcione aunque la instancia de EC2 con MLflow esté detenida.
Durante la verificación se encontró que \texttt{mlflow} sin versión fija
instalaba la 3.16.1 en el contenedor, cuyo \texttt{list\_artifacts()} devolvía
una lista vacía y hacía fallar el arranque; se fijó \texttt{mlflow==3.15.1}.

% =====================================================================
\section{Principales resultados y conclusiones}

\begin{itemize}
\item El desempeño esperable es modesto: ROC-AUC de 0.69 y recall de 0.64 con
precisión de 0.65. Con n=89 y un error estándar del recall cercano a 0.18, las
diferencias entre los mejores modelos no son concluyentes.
\item La conclusión de la Entrega 2 no se sostuvo. Con el holdout de 18 filas,
la regresión logística parecía el mejor baseline y el Random Forest parecía
sobreajustado; con validación cruzada repetida el orden se invierte entre los
baselines. Un único split pequeño no permite comparar modelos.
\item El riesgo lo explican sobre todo el peso al nacer y la duración de la
hospitalización y del oxígeno. Esto es coherente con la literatura y con
predecir al alta, momento en que esas variables ya se conocen.
\item Las probabilidades del modelo quedan entre 0.46 y 0.75 por lo débil de la
señal, así que el puntaje debe leerse de forma relativa (qué paciente está más
arriba en la cohorte) y no como probabilidad absoluta. Esto respalda el uso
del listado de priorización sobre el puntaje aislado.
\item Limitaciones: muestra pequeña de un solo hospital, complicaciones raras
cuyo efecto no se puede estimar con precisión y ausencia de validación externa.
Los siguientes pasos serían una cohorte multicéntrica, calibración de
probabilidades y fijar también la versión de \texttt{scikit-learn} en el
despliegue para que el artefacto se cargue con la misma versión con que se
entrenó.
\end{itemize}

% =====================================================================
\section{Repositorio y evidencias}

\begin{itemize}
\item Repositorio: \url{https://github.com/MAIA-FinalProject/Microproyecto}
(código de datos, features, modelos, API, tablero y artefactos de despliegue).
\item Datos versionados con DVC, con remoto en Google Drive; el pipeline de
features está definido como stage en \texttt{dvc.yaml}.
\item MLflow: servidor en EC2, experimento \texttt{optimizacion\_modelo\_final};
run del modelo final:\\ \texttt{b234c520c8264c999e98fbf545bc226a}.
\item Manual de usuario, manual de instalación y video de síntesis:
\href{https://drive.google.com/file/d/114ZuO4xJnnT2z76jxk2Ibz5pW-Ksr1J5/view}{enlace en Google Drive}.
\end{itemize}

% =====================================================================
\section{Reporte de trabajo en equipo}

El trabajo se organizó en issues de GitHub con revisión cruzada obligatoria
antes de cada merge a \texttt{main}.
\textbf{Martín Quintero} realizó la optimización y selección del modelo final
(tuning, estudio de variables, odds ratios y registro en MLflow, \#21) y
previamente migró el remoto de DVC a Google Drive.
\textbf{Mateo Hernández} desarrolló la integración del tablero con la API, la
explicación local de contribuciones, los nuevos campos del formulario y el
listado de priorización (\#23), además de mantener el repositorio y la CI.
\textbf{María Paula Saa} integró y verificó el despliegue con Docker Compose y
corrigió el bug de versión de MLflow (\#22), además del EDA enfocado.
\textbf{Juan Pablo Buriticá} configuró el servidor de MLflow en EC2 (\#11) y
elaboró los manuales de usuario e instalación y el video de síntesis (\#25).
\textbf{Andrés Romero} consolidó el reporte final a partir del feedback y de los
pull requests (\#24), revisó e integró los PR de modelo final y documentación,
y previamente desarrolló el EDA general y el pipeline de feature engineering.

\begin{thebibliography}{3}

\bibitem{darabi2024}
Darabi, A., Faramarzi, R., Boskabadi, H., Maamouri, G., Rezvani, R.:
Dataset on neonatal and maternal factors influencing neurodevelopmental outcomes
in preterm infants. Data in Brief \textbf{53}, 110058 (2024).
\doi{10.1016/j.dib.2024.110058}

\bibitem{faramarzi2023}
Faramarzi, R., Darabi, A., Emadzadeh, M., Maamouri, G., Rezvani, R.:
Predicting neurodevelopmental outcomes in preterm infants: a comprehensive
evaluation of neonatal and maternal risk factors. Early Human Development
\textbf{184}, 105834 (2023).
\doi{10.1016/j.earlhumdev.2023.105834}

\bibitem{bayley}
Bayley, N.: Bayley Scales of Infant and Toddler Development, 3rd edn.

\end{thebibliography}

\end{document}